\documentclass[11pt,a4paper]{article}
\usepackage{amsmath,amssymb,amsfonts}
\usepackage{graphicx}
\usepackage[margin=1in]{geometry}

\title{Testing the Integer Crossing Hypothesis:\\Does $\alpha^{-1}(\mu) = 137$ at Any Energy Scale?}

\author{S-wind Conjecture Verification System\\Valhalla Project}

\date{\today}

\begin{document}

\maketitle

\begin{abstract}
We test the \textbf{Integer Crossing Hypothesis}: the conjecture that there exists an energy scale $\mu^*$ such that the running fine structure constant inverse $\alpha^{-1}(\mu^*) = 137$ exactly. This hypothesis is motivated by the algebraic fact that $D^4 - 2D^3 + 2D + 1 = 137$ at $D=4$ and the empirical observation that $\alpha^{-1}(m_e) = 137.036$ (0.026\% from 137). Using Standard Model renormalization group running, we demonstrate that this hypothesis is \textbf{falsified}: $\alpha^{-1}(\mu)$ monotonically decreases from 137.036 at $\mu = m_e$ and never crosses 137 at any physical scale. We also rule out cosmological time variation as an explanation, as the required drift exceeds observational bounds by a factor of $\sim 26$. The proximity of $\alpha^{-1}$ to the integer 137 remains an unexplained coincidence.
\end{abstract}

%==============================================================================
\section{Introduction}
%==============================================================================

The fine structure constant $\alpha \approx 1/137$ has fascinated physicists since its identification as a fundamental dimensionless parameter of electromagnetism. The fact that $\alpha^{-1} \approx 137.036$ is tantalizingly close to the integer 137 has motivated numerous attempts to derive this value from first principles.

In previous work, we established the algebraic identity:
\begin{equation}
D^4 - 2D^3 + 2D + 1 = 137 \quad \text{(exactly, for } D=4\text{)}
\end{equation}

This raises a natural question: Is there a physical mechanism that would predict $\alpha^{-1} = 137$ exactly? One possibility is that the running coupling $\alpha^{-1}(\mu)$ crosses 137 at some special energy scale $\mu^*$.

%==============================================================================
\section{The Integer Crossing Hypothesis}
%==============================================================================

\textbf{Conjecture (Integer Crossing Hypothesis):}
\textit{There exists an energy scale $\mu^*$ such that}
\begin{equation}
\alpha^{-1}(\mu^*) = 137 \quad \text{(exactly)}
\end{equation}

\textbf{Motivation:}
\begin{itemize}
    \item $D^4 - 2D^3 + 2D + 1 = 137$ (algebraic fact, $D=4$)
    \item $\alpha^{-1}(m_e) = 137.035999084 \pm 0.000000021$ (CODATA 2022)
    \item Deviation from 137: only 0.026\%
\end{itemize}

%==============================================================================
\section{QED Running Coupling}
%==============================================================================

In quantum electrodynamics, the fine structure constant runs with energy scale $\mu$ according to the renormalization group equation. At one-loop:
\begin{equation}
\alpha^{-1}(\mu) = \alpha^{-1}(\mu_0) - \frac{\beta_0}{2\pi} \ln\left(\frac{\mu}{\mu_0}\right)
\end{equation}
where
\begin{equation}
\beta_0 = \frac{4}{3} \sum_f Q_f^2 N_c
\end{equation}
is the QED beta function coefficient, summed over active fermions with electric charge $Q_f$ and color multiplicity $N_c$.

\subsection{Key Physical Features}

\begin{enumerate}
    \item \textbf{Below $m_e$}: No charged fermions are active, so $\beta_0 \approx 0$. The coupling is essentially constant.
    
    \item \textbf{Above $m_e$}: Electron loops contribute $\beta_0 = 4/3$, causing $\alpha^{-1}$ to \textit{decrease} with increasing $\mu$.
    
    \item \textbf{Direction}: $\alpha^{-1}(\mu)$ monotonically decreases as $\mu$ increases due to vacuum polarization screening.
\end{enumerate}

\subsection{Measured Values}

\begin{center}
\begin{tabular}{|l|c|c|}
\hline
Scale & $\mu$ & $\alpha^{-1}(\mu)$ \\
\hline
Electron mass & $m_e = 0.511$ MeV & 137.036 \\
Z boson mass & $M_Z = 91.2$ GeV & 127.95 \\
\hline
\end{tabular}
\end{center}

%==============================================================================
\section{Analysis: Where Does $\alpha^{-1} = 137$?}
%==============================================================================

\subsection{High Energy Direction}

Since $\alpha^{-1}$ decreases with increasing $\mu$, and $\alpha^{-1}(m_e) = 137.036 > 137$, there is no scale $\mu > m_e$ where $\alpha^{-1} = 137$.

\subsection{Low Energy Direction}

To find $\alpha^{-1} = 137$, we would need $\mu < m_e$. However, below the electron mass threshold, the electron decouples from the running:
\begin{equation}
\beta_0 \to 0 \quad \text{for } \mu < m_e
\end{equation}

This means $\alpha^{-1}$ is essentially \textbf{constant} below $m_e$, frozen at approximately 137.036.

\subsection{Conclusion}

\begin{equation}
\boxed{\text{There is no scale } \mu \text{ where } \alpha^{-1}(\mu) = 137 \text{ in standard QED.}}
\end{equation}

%==============================================================================
\section{Full Standard Model Running}
%==============================================================================

We verify this conclusion using the full Standard Model running with all particle thresholds:

\begin{center}
\begin{tabular}{|l|c|c|}
\hline
Scale & $\mu$ (GeV) & $\alpha^{-1}(\mu)$ \\
\hline
$m_e$ & 0.000511 & 137.036 \\
$m_\mu$ & 0.106 & 135.90 \\
$m_\tau$ & 1.78 & 133.47 \\
$m_b$ & 4.18 & 132.30 \\
$M_W$ & 80.4 & 128.07 \\
$M_Z$ & 91.2 & 127.95 \\
$m_t$ & 173 & 127.50 \\
\hline
\end{tabular}
\end{center}

\textbf{Observation}: $\alpha^{-1}(\mu)$ monotonically decreases from 137.036. It never equals 137 at any physical scale.

%==============================================================================
\section{Alternative: Cosmological Time Variation}
%==============================================================================

An alternative hypothesis is that $\alpha^{-1}$ \textit{was} exactly 137 at some cosmological epoch and has drifted to 137.036 over cosmic time.

\subsection{Required Drift}

\begin{equation}
\frac{\Delta\alpha}{\alpha} = \frac{137.036 - 137}{137} = 2.63 \times 10^{-4}
\end{equation}

\subsection{Observational Constraints}

Measurements from quasar absorption spectra constrain time variation of $\alpha$ over cosmological timescales:
\begin{equation}
\left|\frac{\Delta\alpha}{\alpha}\right| < 10^{-5} \quad \text{(at redshift } z \sim 0.5-3.5\text{)}
\end{equation}

\subsection{Conclusion}

The required drift ($2.63 \times 10^{-4}$) exceeds the observational bound ($10^{-5}$) by a factor of $\sim 26$.

\begin{equation}
\boxed{\text{Cosmological drift hypothesis is RULED OUT by current observations.}}
\end{equation}

%==============================================================================
\section{Summary of Results}
%==============================================================================

\subsection{Hypothesis Tested}

\textbf{Integer Crossing Hypothesis}: $\exists \mu^*$ such that $\alpha^{-1}(\mu^*) = 137$ exactly.

\subsection{Result}

\begin{center}
\textbf{FALSIFIED within the Standard Model}
\end{center}

\subsection{Reasons}

\begin{enumerate}
    \item $\alpha^{-1}(\mu)$ runs from 137.036 at $m_e$ \textit{downward} to $\sim 128$ at $M_Z$
    \item Below $m_e$, running stops ($\beta_0 \to 0$), freezing $\alpha^{-1} \approx 137.036$
    \item No energy scale exists where $\alpha^{-1} = 137$ exactly
    \item Cosmological time variation is ruled out by quasar observations
\end{enumerate}

%==============================================================================
\section{What Remains Unexplained}
%==============================================================================

\subsection{Established Facts}

\begin{enumerate}
    \item $D^4 - 2D^3 + 2D + 1 = 137$ (algebraic identity, $D=4$)
    \item $\alpha^{-1}(m_e) = 137.036$ (measured)
    \item Deviation: 0.026\%
\end{enumerate}

\subsection{Open Questions}

\begin{enumerate}
    \item Why is $\alpha^{-1}$ so close to the integer 137?
    \item Is there a deeper connection to spacetime dimension $D=4$?
    \item What physics beyond the Standard Model could predict $\alpha^{-1} = 137$?
\end{enumerate}

%==============================================================================
\section{Alternative Conjectures for Future Work}
%==============================================================================

\subsection{Conjecture A: UV Integer Quantization}

The ``bare'' coupling at the UV cutoff (Planck scale) might be exactly 137, with radiative corrections producing the observed 137.036.

\textbf{Test}: Compute $\alpha^{-1}(M_{Pl})$ in various BSM scenarios.

\subsection{Conjecture B: Dimension-Coupling Connection}

The $D=4$ polynomial identity suggests a deep connection between spacetime dimension and electromagnetic coupling.

\textbf{Test}: Derive a theoretical mechanism linking $D$ and $\alpha$.

\subsection{Conjecture C: Anthropic Selection}

$\alpha^{-1}$ being near 137 might reflect anthropic constraints on chemistry and nuclear physics.

\textbf{Test}: Compute sensitivity of complex chemistry to $\alpha$ variations.

%==============================================================================
\section{Conclusion}
%==============================================================================

We have tested the Integer Crossing Hypothesis using Standard Model physics and found it to be \textbf{falsified}. The running coupling $\alpha^{-1}(\mu)$ never equals 137 at any physical energy scale, and cosmological time variation cannot account for the 0.026\% deviation.

However, the proximity of $\alpha^{-1}$ to the integer 137 remains unexplained. The algebraic connection $D^4 - 2D^3 + 2D + 1 = 137$ at $D=4$ is a tantalizing hint that may point toward new physics.

This work demonstrates the value of \textbf{systematic conjecture testing}: even a falsified hypothesis sharpens our understanding and points toward new directions.

\begin{center}
\textit{Conjecture $\to$ Test $\to$ Falsify $\to$ Refine $\to$ Repeat}
\end{center}

\section*{Acknowledgments}

Generated by S-wind v1.1.1 Conjecture Verification System.

\end{document}
